\documentclass{beamer}
\usepackage[T1]{fontenc}
\usepackage[utf8]{inputenc}
\usepackage{lmodern}
\usepackage{amsmath,amssymb,mathtools,bm}
\usepackage{booktabs,array,listings,xcolor}
\usetheme{Boadilla}
\useoutertheme{miniframes}
\setbeamertemplate{navigation symbols}{}
\setbeamertemplate{itemize items}[default]
\definecolor{codeblue}{RGB}{42,55,150}
\lstset{language=Python,basicstyle=\ttfamily\scriptsize,keywordstyle=\color{codeblue},
  showstringspaces=false,columns=fullflexible,keepspaces=true,frame=single,
  breaklines=true,aboveskip=3pt,belowskip=2pt}
\newcommand{\BS}{\mathcal L_{\!BS}}
\newcommand{\R}{\mathbb R}
\newcommand{\dd}{\mathrm d}
\newcommand{\pos}[1]{(#1)^+}
\newcommand{\Qtitle}[1]{\textbf{#1}}
\title[MATH4907 Tutorial 4]{MATH4907 Tutorial 4}
\author{Sanyang Liu}
\date{November 6, 2026}
\AtBeginSection[]{\begin{frame}{Roadmap}\tableofcontents[currentsection,hideallsubsections]\end{frame}}

\begin{document}
\maketitle

\section{Q1: American put}

\begin{frame}{Q1(a): Delta hedge -- Problem}
\small
An American put with strike $K$ and expiry $T$ has value $V(S,t)$.
The non-dividend-paying stock price process $S$ satisfies
\[
  \dd S=\mu S\,\dd t+\sigma S\,\dd W,
\]
where $W$ is Brownian motion. In $V(S,t)$, $V_t=\partial V/\partial t$
denotes the time derivative at fixed $S$.
\begin{block}{Problem \hfill 3 points}
Construct the hedge $\Pi=V(S,t)-\Delta S$ with $\Delta=V_S(S,t)$.
Use It\^o's formula to derive
\[
  V_t+\tfrac12\sigma^2 S^2 V_{SS}+rSV_S-rV\leq 0.
\]
Why can this inequality be strict for an American option?
\end{block}
\end{frame}

\begin{frame}{Q1(a): Delta hedge -- Solution}
\small
Along the stock-price path, It\^o's formula gives
\[
\dd V=\bigl(V_t+\mu S V_S+\tfrac12\sigma^2S^2V_{SS}\bigr)\dd t
  +\sigma S V_S\,\dd W.
\]
Over the step, hold one option and $-\Delta$ shares, with
$\Delta=V_S(S,t)$ evaluated at the start of the step. The hedged gain is
\[
\dd\Pi=\dd V-\Delta\,\dd S
       =\bigl(V_t+\tfrac12\sigma^2S^2V_{SS}\bigr)\dd t.
\]
The Brownian term and the physical drift $\mu$ both cancel.
\end{frame}

\begin{frame}{Q1(a): Why the pricing relation is an inequality}
\small
The hedged gain cannot exceed the risk-free return on its value; otherwise
buying the hedge with borrowed funds gives an arbitrage. Thus
\[
V_t+\tfrac12\sigma^2S^2V_{SS}
 \leq r(V-SV_S),
\qquad
\boxed{\BS V:=V_t+\tfrac12\sigma^2S^2V_{SS}+rSV_S-rV\leq0.}
\]
The reverse arbitrage argument for a European option does not apply everywhere:
the American holder controls exercise time. In the continuation region,
$\BS V=0$. In the interior exercise region, $V=K-S$, so
\[
\BS(K-S)=-rK<0\qquad (r>0).
\]
\end{frame}

\begin{frame}{Q1(b): Lower bounds for the American put}
\small
Write $p(S,t)$ for the corresponding European put price.
\begin{block}{Problem \hfill 2 points}
Use no-arbitrage to explain both $V(S,t)\geq\pos{K-S}$ and $V(S,t)\geq p(S,t)$.
\end{block}
\begin{itemize}
\item $V(S,t)\geq\pos{K-S}$ because the holder can exercise immediately.
\item $V(S,t)\geq p(S,t)$ because the holder can wait until expiry and
receive the same payoff as the European put.
\end{itemize}
\end{frame}

\begin{frame}{Q1(c): Exercise and continuation regions}
\small
\begin{block}{Problem \hfill 3 points}
Define the optimal exercise boundary
\[
  S_f(t)=\sup\{S\in[0,K]:V(S,t)=\pos{K-S}\}.
\]
Using $S_f(t)$, define the \textbf{exercise region} and
\textbf{continuation region}; explain what the holder does in each region.
\end{block}
For $t<T$:
\begin{itemize}
\item \textbf{Exercise region} ($0\leq S\leq S_f(t)$): immediate exercise
is optimal and $V(S,t)=K-S$.
\item \textbf{Continuation region} ($S>S_f(t)$): keep the option alive
since $V(S,t)>\pos{K-S}$.
\end{itemize}
\end{frame}

\begin{frame}{Q1(d): Free-boundary problem}
\small
\begin{block}{Problem \hfill 4 points}
State the continuation PDE, value matching, terminal payoff and smooth
pasting for the American put.
\end{block}
Determine both $V(S,t)$ and the unknown curve $S_f(t)$:
\begin{align*}
\BS V&=0 && S>S_f(t),\quad 0\leq t<T,\\
V(S_f(t),t)&=K-S_f(t) &&\text{value matching},\\
V_S(S_f(t),t)&=-1 &&\text{smooth pasting},\\
V(S,T)&=\pos{K-S} && S\geq0.
\end{align*}
For the no-dividend put, $S_f(T)=K$ in the terminal-limit convention.
Value matching and smooth pasting apply for $t<T$; the terminal payoff
has a kink at $S=K$. In the exercise region, $V=K-S$ and
$\BS V=-rK$ when $r>0$.
\end{frame}

\section{Q2: Front fixing}

\begin{frame}{Q2(a): Landau transformation -- Problem}
\small
Write $\beta(t)=S_f(t)$ for the unknown exercise boundary.
Use the Landau transformation
\[
  x=\frac{S}{\beta(t)},\qquad S=x\beta(t),\qquad
  u(x,t)=V(x\beta(t),t).
\]
The moving continuation region $S>\beta(t)$ becomes $x>1$.
\begin{block}{Problem \hfill 3 points}
By the chain rule, show that
\[
 xu_x=SV_S,\qquad x^2u_{xx}=S^2V_{SS},\qquad
 u_t-\frac{x\beta'(t)}{\beta(t)}u_x=V_t.
\]
\end{block}
\end{frame}

\begin{frame}{Q2(a): Chain-rule identities}
\small
Differentiate while keeping the other transformed variable fixed:
\begin{align*}
u_x&=\beta V_S,& xu_x&=SV_S,\\
u_{xx}&=\beta^2V_{SS},&x^2u_{xx}&=S^2V_{SS},\\
u_t&=V_t+x\beta'V_S
    =V_t+\frac{x\beta'}{\beta}u_x.
\end{align*}
Therefore
\[
\boxed{V_t=u_t-\frac{x\beta'}{\beta}u_x.}
\]
$V_t$ is taken at fixed $S$, and $u_t$ at fixed $x$.
\end{frame}

\begin{frame}{Q2(b): PDE on the fixed domain}
\small
\begin{block}{Problem \hfill 3 points}
Use the identities in part (a) to derive the transformed PDE and state its domain.
\end{block}
Substitute the three identities into $\BS V=0$:
\[
\boxed{
u_t+x\left(r-\frac{\beta'}{\beta}\right)u_x
       +\frac{\sigma^2x^2}{2}u_{xx}-ru=0,
\qquad 1<x<\infty,\quad 0<t<T.}
\]
The $-x\beta' u_x/\beta$ term is due to the moving coordinate.
Omitting it would solve a different pricing problem.
\end{frame}

\begin{frame}{Q2(c): Transformed conditions}
\small
\begin{block}{Problem \hfill 3 points}
Transform the value matching, smooth pasting and terminal payoff from Q1(d)
into conditions on $u(x,t)$ and $\beta(t)$.
\end{block}
At $S=\beta(t)$, we have $x=1$. Consequently
\[
\boxed{u(1,t)=K-\beta(t)},\qquad
\boxed{u_x(1,t)=-\beta(t)},\qquad t<T.
\]
At expiry $\beta(T)=K$, hence on the continuation strip $x\geq1$,
\[
u(x,T)=\pos{K-x\beta(T)}=\pos{K(1-x)}=0.
\]
Smooth pasting is a pre-expiry condition; the payoff is not smooth at the
corner $(x,t)=(1,T)$.
\end{frame}

\begin{frame}{Q2(d): Finite-difference stencils -- Problem}
\small
\begin{block}{Problem \hfill 4 points}
Introduce an artificial boundary $x=x_r$ with $u(x_r,t)=0$.
Define
\[
 h=\frac{x_r-1}{n+1},\quad k=\frac{T}{m},\qquad
 x_i=1+ih,\;t_j=jk,
\]
where $i=0,\ldots,n+1$ and $j=0,\ldots,m$. Write down
finite-difference approximations for $u_x$, $u_{xx}$ and $u_t$
suitable for discretizing the transformed PDE.
\end{block}
The put value approaches zero as the stock price becomes large.
\end{frame}

\begin{frame}{Q2(d): Finite-difference stencils}
\small
Let $x_i=1+ih$, $t_j=jk$, $h=(x_r-1)/(n+1)$, $k=T/m$, and
$U_{i,j}\approx u(x_i,t_j)$. Centered spatial differences give
\[
u_x(x_i,t_j)\approx\frac{U_{i+1,j}-U_{i-1,j}}{2h},\qquad
u_{xx}(x_i,t_j)\approx\frac{U_{i-1,j}-2U_{i,j}+U_{i+1,j}}{h^2}.
\]
Because the payoff is known at $t=T$, backward stepping in calendar time
can use
\[
u_t(x_i,t_{j+1})\approx\frac{U_{i,j+1}-U_{i,j}}{k},
\qquad j=m-1,m-2,\ldots,0.
\]
\end{frame}

\begin{frame}{Q2(d): Coupled discrete equations}
\footnotesize
One explicit backward-time discretization is
\begin{align*}
0={}&\frac{U_{i,j+1}-U_{i,j}}{k}
+x_i\left(r-\frac{\beta_{j+1}-\beta_j}{k\beta_{j+1}}\right)
\frac{U_{i+1,j+1}-U_{i-1,j+1}}{2h}\\
&+\frac{\sigma^2x_i^2}{2}
\frac{U_{i-1,j+1}-2U_{i,j+1}+U_{i+1,j+1}}{h^2}
-rU_{i,j+1}.
\end{align*}
Use $U_{0,j}=K-\beta_j$, $U_{n+1,j}=0$, and the one-sided smooth-pasting
condition
\[
\frac{U_{1,j}-U_{0,j}}{h}\approx-\beta_j,\qquad j<m.
\]
These equations determine $\beta_j$ and the new interior values together.
\end{frame}

\section{Q3: Complementarity}

\begin{frame}{Q3(a): Heat obstacle -- Problem}
\small
Apply the change of variables
\[
 x=\ln(S/K),\quad \tau=\tfrac12\sigma^2(T-t),\quad
 u=e^{-\alpha x-\beta\tau}V(S,t),
\]
where $\alpha=\tfrac12-r/\sigma^2$ and
$\beta=-(\tfrac12+r/\sigma^2)^2$. The American put satisfies
\[
 u_\tau-u_{xx}\geq0,\qquad u-g\geq0,\qquad
 (u_\tau-u_{xx})(u-g)=0,
\]
with $g(x,\tau)=Ke^{-\alpha x-\beta\tau}\pos{1-e^x}$.
\begin{block}{Problem \hfill 3 points}
Explain the financial meaning of $u_\tau-u_{xx}\geq0$ and $u-g\geq0$.
Relate the two inequalities to the exercise and continuation regions.
\end{block}
\end{frame}

\begin{frame}{Q3(a): Financial meaning of complementarity}
\small
The change of variables multiplies the obstacle inequality by a
positive factor and transforms the pricing operator according to
\[
\BS V=-\frac{\sigma^2}{2}e^{\alpha x+\beta\tau}(u_\tau-u_{xx}).
\]
\begin{itemize}
\item $u-g\geq0$: the option cannot be worth less than immediate exercise.
\item $u_\tau-u_{xx}\geq0$: equivalent to $\BS V\leq0$.
\item \textbf{Continue:} $u>g$ forces $u_\tau-u_{xx}=0$.
\item \textbf{Exercise:} $u=g$ allows a positive heat residual.
\end{itemize}
For $V=K-S$, $\BS V=-rK$, and hence the heat residual is strictly
positive when $r>0$.
\end{frame}

\begin{frame}{Q3(b): Artificial log-price boundaries}
\small
\begin{block}{Problem \hfill 3 points}
For artificial boundaries $x_l<0<x_r$, explain why it is reasonable
to impose $u(x_r,\tau)=0$ and $u(x_l,\tau)=g(x_l,\tau)$.
\end{block}
Choose $x_l<0<x_r$. Since $S=Ke^x$:
\[
\boxed{u(x_r,\tau)=0}
\]
approximates the small value of a put at large stock prices. At a sufficiently
negative $x_l$, the stock is deep in the exercise region, so
\[
\boxed{u(x_l,\tau)=g(x_l,\tau)}.
\]
These conditions truncate an infinite spatial domain. They introduce
boundary error if $x_l$ or $x_r$ is too close to the option's relevant region.
\end{frame}

\begin{frame}{Q3(c): Backward Euler in transformed time}
\small
\begin{block}{Problem \hfill 4 points}
Use a backward difference in $\tau$ and a centered difference in $x$
to derive the discrete inequality for $u_\tau-u_{xx}\geq0$.
\end{block}
With $\lambda=k/h^2$, apply backward Euler at $(x_i,\tau_j)$:
\[
\frac{u_{i,j}-u_{i,j-1}}{k}
-\frac{u_{i-1,j}-2u_{i,j}+u_{i+1,j}}{h^2}\geq0.
\]
Multiplication by $k$ yields
\[
\boxed{-\lambda u_{i-1,j}+(1+2\lambda)u_{i,j}
-\lambda u_{i+1,j}\geq u_{i,j-1}.}
\]
The obstacle is imposed at the new time: $u_{i,j}\geq g(x_i,\tau_j)$.
\end{frame}

\begin{frame}{Q3(d): Matrix LCP}
\small
\begin{block}{Problem \hfill 3 points}
Write the discrete problem for the interior vector $\bm u^j$ as a matrix
complementarity system, including the left boundary term. Define $A$.
\end{block}
Let $\bm u^j=(u_{1,j},\ldots,u_{n,j})^\top$ and
$\bm g^j=(g(x_1,\tau_j),\ldots,g(x_n,\tau_j))^\top$. Define
\[
A=\begin{pmatrix}
1+2\lambda&-\lambda&&\\-\lambda&1+2\lambda&\ddots&\\
&\ddots&\ddots&-\lambda\\&&-\lambda&1+2\lambda
\end{pmatrix},\quad
\bm b^j=\bm u^{j-1}+\lambda g(x_l,\tau_j)\bm e_1.
\]
The right artificial boundary contributes nothing because $u(x_r,\tau_j)=0$.
The discrete problem is
\[
A\bm u^j-\bm b^j\geq0,\quad \bm u^j-\bm g^j\geq0,\quad
(A\bm u^j-\bm b^j)^\top(\bm u^j-\bm g^j)=0.
\]
\end{frame}

\begin{frame}{Q3(e): Shift to a nonnegative unknown}
\small
\begin{block}{Problem \hfill 2 points}
Set $\bm z^j=\bm u^j-\bm g^j$. Rewrite the problem in standard LCP form and
give the vector $\bm q^j$ explicitly.
\end{block}
Let $\bm z^j=\bm u^j-\bm g^j\geq0$. Substitute
$\bm u^j=\bm z^j+\bm g^j$ into the matrix residual:
\[
A\bm u^j-\bm b^j=A\bm z^j+\bm q^j,
\]
where
\[
\boxed{\bm q^j=A\bm g^j-\bm u^{j-1}
-\lambda g(x_l,\tau_j)\bm e_1.}
\]
The standard linear complementarity form is
\[
\bm q^j+A\bm z^j\geq0,\qquad \bm z^j\geq0,\qquad
(\bm q^j+A\bm z^j)^\top\bm z^j=0.
\]
At each component, at least one of the two nonnegative quantities is zero.
\end{frame}

\section{Q4: Projected methods}

\begin{frame}{Q4(a): Scalar complementarity -- Problem}
\small
The general linear complementarity problem (LCP) is
\[
 \bm p+M\bm z\geq0,\quad \bm z\geq0,\quad
 (\bm p+M\bm z)^\top\bm z=0,
\]
equivalently $\min\{\bm z,\bm p+M\bm z\}=\bm0$ componentwise.
\begin{block}{Problem \hfill 2 points}
For the scalar case $n=1$, $M=[m]$ with $m>0$, solve
$\min\{z,p+mz\}=0$ and prove
\[
 z=\max\{0,-p/m\}.
\]
\end{block}
\end{frame}

\begin{frame}{Q4(a): One-dimensional complementarity}
\small
For $m>0$, solve $\min\{z,p+mz\}=0$. Complementarity says
\[
z\geq0,\quad p+mz\geq0,\quad z(p+mz)=0.
\]
If $p\geq0$, choose $z=0$. If $p<0$, then $z=0$ violates the second
inequality, so $p+mz=0$ and $z=-p/m>0$. Therefore
\[
\boxed{z=\max\{0,-p/m\}.}
\]
The maximum is the projection of the unconstrained root onto $[0,\infty)$.
\end{frame}

\begin{frame}{Q4(b): Triangular example}
\small
\begin{block}{Problem \hfill 3 points}
Solve $\min\{\bm z,\bm p+M\bm z\}=\bm0$ for
\[
 M=\begin{pmatrix}3&0&0\\1&3&0\\1&2&3\end{pmatrix},\qquad
 \bm p=\begin{pmatrix}-3\\2\\-1\end{pmatrix}.
\]
\end{block}
Solve by forward componentwise projection:
\begin{align*}
z_1&=\max\{0,-(-3)/3\}=1,\\
z_2&=\max\{0,-(2+z_1)/3\}=0,\\
z_3&=\max\{0,-(-1+z_1+2z_2)/3\}=0.
\end{align*}
Thus $\boxed{\bm z=(1,0,0)^\top}$, and
$\bm p+M\bm z=(0,3,0)^\top\geq0$ verifies complementarity.
\end{frame}

\begin{frame}{Q4(c): Projected Jacobi}
\small
\begin{block}{Problem \hfill 2 points}
Write down the \textbf{projected Jacobi iteration} for the general LCP.
Specify which iterate is used in the off-diagonal terms and how the
nonnegativity constraint is enforced.
\end{block}
Split $M=L+D+U$ into its strictly lower, diagonal and strictly upper parts,
with $m_{ii}>0$.
Using the old iterate for all off-diagonal terms gives
\[
\boxed{z_i^{(\ell+1)}=
\max\left\{0,-\frac{p_i+\sum_{j\ne i}m_{ij}z_j^{(\ell)}}{m_{ii}}\right\},
\quad i=1,\ldots,n.}
\]
Equivalently, each component is projected onto the nonnegative half-line
after the Jacobi update. Start with $\bm z^{(0)}\geq0$.
\end{frame}

\begin{frame}{Q4(d): Projected Gauss--Seidel}
\small
\begin{block}{Problem \hfill 3 points}
Write down the \textbf{projected Gauss--Seidel iteration}.
Specify when new vs. old components are used. Explain why the method
is called ``projected.''
\end{block}
Use newly computed components to the left and old ones to the right:
\[
\boxed{z_i^{(\ell+1)}=
\max\left\{0,-\frac{p_i+
\sum_{j<i}m_{ij}z_j^{(\ell+1)}+
\sum_{j>i}m_{ij}z_j^{(\ell)}}{m_{ii}}\right\}.}
\]
It is \emph{projected} because $\max\{0,\cdot\}$ enforces
$z_i^{(\ell+1)}\geq0$ after each scalar update.
For a triangular $M$, one Gauss--Seidel sweep from top to bottom solves
the example in part (b).
\end{frame}

\section{Q5: Barrier call}

\begin{frame}{Q5(a): Down-and-out call problem}
\small
\begin{block}{Problem \hfill 3 points}
A European down-and-out call with strike $K$, barrier $H$ and expiry $T$
is extinguished when $S(t)$ reaches $H$ from above. State the Black--Scholes
PDE, the terminal payoff and the condition at the barrier.
\end{block}
Assume $H<K$ (as in the numerical data). The option survives while $S>H$:
\begin{align*}
V_t+\tfrac12\sigma^2S^2V_{SS}+rSV_S-rV&=0,
&&S>H,\quad 0<t<T,\\
V(S,T)&=\pos{S-K},&&S>H,\\
V(H,t)&=0,&&0\leq t\leq T.
\end{align*}
\end{frame}

\begin{frame}{Q5(b): Heat-equation transformation}
\small
\begin{block}{Problem \hfill 3 points}
Use the change of variables below to derive the heat equation and the boundary condition at $x=0$.
\end{block}
Take
\[
x=\log(S/H),\quad \tau=\frac{\sigma^2}{2}(T-t),\quad
V=e^{a\tau+bx}u,\quad
b=\frac12-\frac r{\sigma^2},\quad
a=-b^2-\frac{2r}{\sigma^2}.
\]
After changing variables, the Black--Scholes PDE divided by
$\sigma^2e^{a\tau+bx}/2$ becomes
\[
-u_\tau+u_{xx}+(2b+2r/\sigma^2-1)u_x
 +(b^2+(2r/\sigma^2-1)b-2r/\sigma^2-a)u=0.
\]
The chosen $a,b$ cancel the first- and zeroth-order terms, leaving
$\boxed{u_\tau-u_{xx}=0}$ for $x>0$.
\end{frame}

\begin{frame}{Q5(b): Initial and barrier data}
\small
At $\tau=0$, calendar time is $T$, so
\[
\boxed{u(x,0)=e^{-bx}\pos{He^x-K}},\qquad x>0.
\]
The barrier $S=H$ maps to $x=0$. Since the exponential factor never
vanishes,
\[
\boxed{u(0,\tau)=0},\qquad 0\leq\tau\leq\tau_*:=\frac{\sigma^2T}{2}.
\]
At the requested price $S_0$, evaluate the heat solution at
$x_\star=\log(S_0/H)$ and $\tau_*$, then undo the transformation:
\[
V(S_0,0)=e^{a\tau_*+bx_\star}u(x_\star,\tau_*).
\]
\end{frame}

\begin{frame}{Q5(c): Numerical pricing -- Problem}
\small
\begin{block}{Problem \hfill 5 points}
Implement an \textbf{explicit FDM} for the transformed heat equation with
\[
 S_0:=S(0)=100,\quad K=95,\quad H=80,\quad T=1,\quad
 r=0.05,\quad\sigma=0.30.
\]
Use the artificial boundary
\[
 x_r=(r-\tfrac12\sigma^2)T+8\sigma\sqrt{T},
\]
with a suitable boundary condition at $x=x_r$, and choose the mesh so that
\[
 \lambda=k/h^2\leq\tfrac12.
\]
Compute $V(S_0,0)$ and compare with the \textbf{analytic benchmark}
$15.429266$. Submit the code and briefly explain interpolation when
$S_0$ is not located exactly at a spatial grid point.
\end{block}
\end{frame}

\begin{frame}{Q5(c): Far-field boundary and parameters}
\small
The assignment specifies
\[
x_r=(r-\tfrac12\sigma^2)T+8\sigma\sqrt T=2.405,
\quad S_r=He^{x_r}\simeq886.2744.
\]
At such a large stock value, use $V(S_r,t)\simeq S_r-Ke^{-r(T-t)}$.
Since $T-t=2\tau/\sigma^2$, the transformed boundary is
\[
\boxed{u(x_r,\tau)=e^{-a\tau-bx_r}
\left(He^{x_r}-Ke^{-2r\tau/\sigma^2}\right).}
\]
For the data, $\tau_*=0.045$, $x_\star=\log(1.25)\simeq0.223144$.
\medskip

\footnotesize Lecture 4, slide 54, writes $e^{-r\tau}$ in its last formula.
The stated time scaling requires $e^{-2r\tau/\sigma^2}$.
\end{frame}

\begin{frame}{Q5(c): Explicit heat scheme}
\small
Set $h=x_r/N$, $k=\tau_*/m$, $\lambda=k/h^2$, and
$U_i^j\approx u(ih,jk)$. Then
\[
\boxed{U_i^{j+1}=(1-2\lambda)U_i^j
+\lambda U_{i-1}^j+\lambda U_{i+1}^j},\quad 1\leq i<N.
\]
Initialize $U_i^0=e^{-bih}\pos{He^{ih}-K}$. At each step set
$U_0^{j+1}=0$ and $U_N^{j+1}=u(x_r,(j+1)k)$.
\medskip
\textbf{Stability:} nonnegative weights require $0\leq\lambda\leq1/2$,
as specified in both the assignment and Lecture 4.
\end{frame}

\begin{frame}{Q5(c): Recovering an off-grid price}
\small
Usually $x_\star=\log(S_0/H)$ lies between grid points $x_i$ and $x_{i+1}$.
With $\theta=(x_\star-x_i)/h$, linearly interpolate the heat solution:
\[
u(x_\star,\tau_*)\approx(1-\theta)U_i^m+\theta U_{i+1}^m.
\]
Then multiply by $e^{a\tau_*+bx_\star}$. Interpolating $u$ in $x$ is
consistent with the grid on which the heat equation was solved.
\medskip
With $N=400$ and $m=3113$, the computed price is $15.429952908$,
compared with the analytic benchmark $15.429265555$.
\end{frame}

\begin{frame}[fragile]{Q5(c): Python code -- Setup}
\begin{lstlisting}
from math import ceil, exp, log, sqrt
import numpy as np

S0, K, H, T = 100., 95., 80., 1.
r, sigma = .05, .30
b = .5 - r/sigma**2
a = -b**2 - 2*r/sigma**2
tau_end = sigma**2*T/2
x_right = (r-sigma**2/2)*T + 8*sigma*sqrt(T)

def upper_boundary(tau):
    remaining = 2*tau/sigma**2
    value = H*exp(x_right) - K*exp(-r*remaining)
    return exp(-a*tau-b*x_right)*value
\end{lstlisting}
Run both code blocks together to reproduce the finite-difference prices.
\end{frame}

\begin{frame}[fragile]{Q5(c): Python code -- Time stepping}
\begin{lstlisting}
def price(N, lambda_target=.4):
    h = x_right/N
    m = ceil(tau_end/(lambda_target*h*h))
    k, lam = tau_end/m, tau_end/m/h**2
    x = np.linspace(0., x_right, N+1)
    u = np.exp(-b*x)*np.maximum(H*np.exp(x)-K, 0.)
    u[0] = 0.
    for j in range(m):
        new = np.empty_like(u)
        new[1:-1] = (1-2*lam)*u[1:-1] + lam*(u[:-2]+u[2:])
        new[0] = 0.
        new[-1] = upper_boundary((j+1)*k)
        u = new
    x_eval = log(S0/H)
    value = exp(a*tau_end+b*x_eval)*np.interp(x_eval, x, u)
    return m, lam, value

for N in (200, 400, 600, 800):
    print(N, price(N))
\end{lstlisting}
\end{frame}

\begin{frame}{Q5(d): Why the barrier call is cheaper}
\small
\begin{block}{Problem \hfill 2 points}
Explain why the down-and-out call is cheaper than the corresponding
vanilla European call with the same strike and expiry.
\end{block}
Path by path, the down-and-out payoff is
\[
\pos{S(T)-K}\,\bm 1_{\{\min_{0\leq s\leq T}S(s)>H\}}
\leq\pos{S(T)-K},
\]
the vanilla European call payoff. Discounting and taking risk-neutral
expectations preserves this order:
\[
V_{\rm DO}(S_0,0)\leq C_{\rm vanilla}(S_0,0).
\]
Numerically, $C_{\rm vanilla}\simeq16.801211384$ and
$V_{\rm DO}\simeq15.429265555$. The difference is the value lost on
paths that hit the barrier before expiry.
\end{frame}

\begin{frame}{Q5(e): Numerical convergence}
\small
\begin{block}{Problem \hfill 2 points}
Repeat the computation in (c) with \textbf{at least two different mesh sizes}.
Briefly discuss the \textbf{stability} and \textbf{convergence} of the
computed option prices.
\end{block}
Explicit scheme with target $\lambda=0.4$; $m$ is rounded up, so actual
$\lambda\leq0.4$.
\begin{center}
\begin{tabular}{rrrrr}
\toprule
$N$ intervals & $m$ steps & $\lambda$ & FD price & Absolute error\\\midrule
200&779&0.399489&15.433241249&0.003975693\\
400&3113&0.399874&15.429952908&0.000687352\\
600&7003&0.399945&15.429627064&0.000361508\\
800&12449&0.399970&15.429450136&0.000184580\\\bottomrule
\end{tabular}
\end{center}
The error decreases as the mesh is refined. All runs satisfy the explicit
stability bound. The errors include the effects of discretization,
interpolation and boundary truncation.
\end{frame}

\begin{frame}{Assignment 4: Submission requirements}
\small
\begin{alertblock}{Submission requirements (p.~6)}
\begin{enumerate}
\item Clear derivations for the theoretical questions;
\item the finite-difference code for Question 5;
\item numerical estimates of $V(S_0,0)$;
\item a short discussion of boundary conditions, stability,
interpolation and accuracy.
\end{enumerate}
\end{alertblock}
\end{frame}

\end{document}
